\documentclass[10pt,a4paper]{article}
\usepackage[utf8]{inputenc}\usepackage[T1]{fontenc}\usepackage{lmodern}
\usepackage[spanish,es-nodecimaldot,es-tabla]{babel}
\usepackage[a4paper,margin=2cm]{geometry}
\usepackage{amsmath,amssymb,textcomp,array,tabularx,longtable,booktabs,enumitem,microtype,fancyhdr,titlesec,graphicx,float,caption}
\usepackage[table]{xcolor}
\usepackage[most]{tcolorbox}
\usepackage[colorlinks=true,linkcolor=azul,urlcolor=azul]{hyperref}
\definecolor{azul}{HTML}{1F4E79}\definecolor{celeste}{HTML}{DCE6F2}\definecolor{rojo}{HTML}{C00000}
\definecolor{rosa}{HTML}{FBE9E9}\definecolor{gris}{HTML}{595959}\definecolor{fondoec}{HTML}{F3F6FA}
\newcommand{\diam}{\ensuremath{\varnothing}}
\graphicspath{{../figs/}}
\captionsetup{font=small,labelfont={bf,color=azul}}
\newcommand{\fig}[3]{\begin{figure}[H]\centering\includegraphics[width=#1\linewidth]{#2}\caption{#3}\end{figure}}
\newcommand{\figdos}[5]{\begin{figure}[H]\centering\includegraphics[width=#1\linewidth]{#2}\hfill\includegraphics[width=#3\linewidth]{#4}\caption{#5}\end{figure}}
\newcolumntype{L}{>{\raggedright\arraybackslash}X}
\newcolumntype{C}{>{\centering\arraybackslash}X}
\newcolumntype{R}{>{\raggedleft\arraybackslash}X}
\newcolumntype{P}[1]{>{\raggedright\arraybackslash}p{#1}}
\setlength{\tabcolsep}{4pt}\renewcommand{\arraystretch}{1.25}\arrayrulecolor{azul!45}
\titleformat{\section}{\large\bfseries\color{azul}}{\thesection.}{0.5em}{}
\titleformat{\subsection}{\normalsize\bfseries\color{azul}}{\thesubsection}{0.5em}{}
\newtcolorbox{ecu}{colback=fondoec,colframe=azul!45,boxrule=0.4pt,arc=0pt,left=6pt,right=6pt,top=2pt,bottom=2pt}
\newcommand{\pend}[1]{\textcolor{rojo}{\bfseries #1}}
\newcommand{\nota}[1]{{\par\small\itshape\color{gris}#1\par}\medskip}
\pagestyle{fancy}\fancyhf{}\fancyfoot[R]{\small\color{gris}PDR-STR-REC-2026-v2B · Equipo Hermes · página \thepage}
\renewcommand{\headrulewidth}{0pt}\setlength{\parindent}{0pt}\setlength{\parskip}{4pt}\setlist{nosep,leftmargin=1.4em}
\begin{document}

\begin{tcolorbox}[colback=celeste,colframe=azul,boxrule=0.6pt,arc=0pt]
{\large\bfseries\color{azul} Estructura v2, opción B regulable: estaciones roscadas para ajustar el CG\par}\smallskip
{\bfseries Ajuste del centro de gravedad con la batería y verificación del choque de 30\,G (prueba de caída) ---
CanSat Argentina Universitario 2026\par}
\end{tcolorbox}
\par{\small\noindent\begin{tabularx}{\linewidth}{P{2.8cm}LP{2.8cm}L}
\rowcolor{azul}\color{white}\bfseries Dato & \color{white}\bfseries Valor & \color{white}\bfseries Dato & \color{white}\bfseries Valor \\
\bfseries Código & PDR-STR-REC-2026-v2B & \bfseries Base & Documento v2 del equipo, diseño B (esqueleto) \\
\bfseries Masa & 531,8\,g (recalculada) & \bfseries Envolvente & \diam\,90 $\times$ 240\,mm \\
\bfseries CG nominal & 110,7\,mm desde nadir & \bfseries Rango del CG & 104,5--112,8\,mm \\
\bfseries Rotura en la caída & 427\,G (mamparo M2) & \bfseries Fecha & 21/11/2026 \\
\end{tabularx}}
\par\medskip
\section*{Resumen}
\begin{itemize}
\item \textbf{Propuesta:} las cuatro varillas M3 son roscadas en toda su longitud y cada mamparo queda apretado entre dos
tuercas nyloc. Girando las tuercas se corre la estación: \textbf{0,5\,mm por vuelta}. Se hacen regulables las dos estaciones
que llevan masa: \textbf{M2 con el módulo de batería} (116,8\,g) y \textbf{M3 con el módulo de
aviónica} (77,3\,g).
\item \textbf{Centro de gravedad:} nominal en 110,7\,mm desde el fondo de nadir. Corriendo los dos
módulos se mueve entre 104,5 y 112,8\,mm. El margen estático de la vela
queda positivo en todo el rango.
\item \textbf{Choque de 30\,G (prueba de caída):} con el CanSat colgado de la brida, la opción B regulable rompe recién a
\textbf{427\,G} (mamparo M2). La opción A (monocasco pegado) rompe a
92\,G en el pegado de M4 a la pared.
\item \textbf{Pendientes que no resuelve este cambio:} la masa (531,8\,g, sobre el máximo de 510\,g),
el ensayo de la lijadora con la carena de PETG y el área real de la vela (sección 5).
\end{itemize}

\section{Las dos opciones de v2 y por qué se sigue con B}
El documento v2 compara dos arquitecturas: \textbf{A}, un monocasco de PETG impreso con los mamparos pegados a la pared, y
\textbf{B}, un esqueleto de cuatro varillas M3 con mamparos y una carena de 2\,mm que se desliza por encima. v2 eligió B por
cualitativo (delaminación, acceso a la electrónica). El FEM de la prueba de caída (sección 4) lo confirma con números.

\figdos{0.49}{v2A_interior_anot.png}{0.49}{v2B_interior_anot.png}{Opciones A (izquierda) y B (derecha) del documento v2, en corte: la pared que mira a la cámara está recortada y el interior queda entero.}
\par{\footnotesize\noindent\begin{tabularx}{\linewidth}{P{3.6cm}LL}
\rowcolor{azul}\color{white}\bfseries Criterio & \color{white}\bfseries A: monocasco pegado & \cellcolor{rojo}\color{white}\bfseries B: esqueleto de varillas \\
\cellcolor{celeste}\bfseries Camino de la carga en la caída & Todo cuelga del pegado de M4 a la pared & \cellcolor{rosa}Las varillas llevan la carga a la brida \\
\cellcolor{celeste}\bfseries Rotura en la caída (FEM) & 92\,G (pegado de M4 a la pared) & \cellcolor{rosa}431\,G (mamparo M4) \\
\cellcolor{celeste}\bfseries Acceso a la electrónica & Cerrado al pegar & \cellcolor{rosa}El esqueleto sale entero al sacar 6 tornillos \\
\cellcolor{celeste}\bfseries Ajuste del CG & No (mamparos pegados) & \cellcolor{rosa}Sí, con estaciones roscadas (esta memoria) \\
\cellcolor{celeste}\bfseries Incertidumbre principal & Adherencia sobre PETG (2--6\,MPa) & \cellcolor{rosa}Aflojamiento de tuercas (nyloc + marca) \\
\end{tabularx}}
\par\medskip
\section{Estaciones roscadas}
\subsection{Concepto}
Las varillas M3 de Al 6061 ya eran roscadas en v2; el cambio es usar la rosca en toda la longitud. Cada mamparo queda
apretado entre dos tuercas nyloc con arandela plana, así que su posición se fija en cualquier punto de la varilla.
\textbf{Por qué tuercas y no mamparos roscados:} un mamparo con agujeros roscados no puede girar sobre cuatro varillas
paralelas a la vez (giraría alrededor del eje del CanSat, no de cada varilla). Con tuercas, cada una gira sola sobre su
varilla y el mamparo baja o sube sin girar.
\begin{itemize}
\item \textbf{Módulo de batería:} M2 se imprime con una cuna en U de la que cuelga el pack 2 $\times$ 18650. Se mueve
entero con M2. Las paredes de la cuna contienen las celdas en los costados (carga lateral de la lijadora).
\item \textbf{Módulo de aviónica:} los rieles de las placas cuelgan solo de M3 (en v2 iban de M2 a M3), así que el módulo
se mueve con M3 sin depender de M2.
\item \textbf{Fijos:} M1 (fondo de nadir, ventana de la cámara y pivotes de los largueros) y M4 (copa del paracaídas y
anclaje de la brida). La baliza baja a $x$ = 6--20\,mm para dar recorrido al módulo de batería.
\item \textbf{Herrajes:} ocho tuercas por varilla (32 en total). Las de los extremos son de acero (anclaje de la brida y
apriete del conjunto); las de las estaciones intermedias, nyloc de Al 7075 (0,2\,g). Agregan unos 5\,g respecto de v2.
\end{itemize}

\fig{0.74}{v2BR_interior_anot.png}{Opción B regulable en la posición nominal (corte). M2 y M3 son estaciones regulables sobre las varillas roscadas; M1 y M4 quedan fijos.}
\fig{0.72}{v2BR_detalle_anot.png}{Detalle de la estación de M2: el mamparo queda apretado entre dos tuercas nyloc y la cuna en U, impresa con M2, lleva el pack. Una vuelta de las tuercas corre el módulo 0,5\,mm.}
\subsection{Recorrido de cada módulo}
\par{\footnotesize\noindent\begin{tabularx}{\linewidth}{P{2.0cm}P{1.1cm}P{1.2cm}P{1.6cm}LL}
\rowcolor{azul}\color{white}\bfseries Módulo & \color{white}\bfseries Masa & \color{white}\bfseries Nominal & \color{white}\bfseries Recorrido & \color{white}\bfseries Límite hacia nadir & \color{white}\bfseries Límite hacia cenit \\
\cellcolor{celeste}\bfseries Batería (M2) & 116,8\,g & 83\,mm & 66--90\,mm & La cuna llega a la baliza ($x$ = 20) & Las tuercas de M2 llegan a los rieles de la aviónica \\
\cellcolor{celeste}\bfseries Aviónica (M3) & 77,3\,g & 173\,mm & 156--177\,mm & Los rieles llegan a M2 (hace falta $x_{M3} \geq x_{M2} + 86$) & Las tuercas de M3 llegan a la cápsula de nicrom \\
\end{tabularx}}
\par\medskip
\subsection{Procedimiento de ajuste}
\begin{enumerate}
\item Integrar todo con las estaciones nominales y cerrar la carena.
\item Medir la masa y el CG (sobre un filo o con dos balanzas).
\item Calcular el corrimiento del módulo:
\begin{ecu}
\begin{equation}
\Delta x_\text{módulo} = \frac{m_\text{total}}{m_\text{módulo}}\,\Delta x_\text{CG}\qquad\Rightarrow\qquad
\text{batería: } \frac{531,8}{116,8} = 4,55\ \text{mm de módulo por mm de CG}
\end{equation}
\end{ecu}
\item Sacar la carena, aflojar las tuercas del módulo y girarlas: 2 vueltas = 1\,mm. Apretar con torque bajo (arandela plana
\diam7 para no marcar el PETG) y marcar la posición con pintura.
\item Volver a medir el CG.
\end{enumerate}

\section{Centro de gravedad}
\par{\footnotesize\begin{longtable}{P{9.6cm}P{1.5cm}P{1.4cm}P{2.0cm}}
\rowcolor{azul}\color{white}\bfseries Ítem & \color{white}\bfseries Masa [g] & \color{white}\bfseries $x$ [mm] & \color{white}\bfseries $m\,x$ [kg\,mm] \\ \endfirsthead
\rowcolor{azul}\color{white}\bfseries Ítem & \color{white}\bfseries Masa [g] & \color{white}\bfseries $x$ [mm] & \color{white}\bfseries $m\,x$ [kg\,mm] \\ \endhead
\cellcolor{celeste}\bfseries Carena PETG 2 mm (240 mm, ranuras y canal) & 163,0 & 118 & 19,23 \\
\cellcolor{celeste}\bfseries M1 fondo de nadir & 13,2 & 2 & 0,03 \\
\cellcolor{celeste}\bfseries M4 mamparo de cenit & 13,2 & 218 & 2,88 \\
\cellcolor{celeste}\bfseries M5 tapa de la copa & 13,2 & 238 & 3,14 \\
\cellcolor{celeste}\bfseries Varillas roscadas M3 Al 6061 (4) & 13,6 & 113 & 1,54 \\
\cellcolor{celeste}\bfseries Tuercas de los extremos (8, acero) & 4,4 & 112 & 0,49 \\
\cellcolor{celeste}\bfseries Tuercas de M1 y M4 (8, Al 7075) & 1,6 & 110 & 0,18 \\
\cellcolor{celeste}\bfseries Paracaídas \diam{}226 + líneas & 8,0 & 229 & 1,83 \\
\cellcolor{celeste}\bfseries Resorte expulsor + horquilla y pasador & 5,0 & 226 & 1,13 \\
\cellcolor{celeste}\bfseries Vela Rogallo plegada & 11,5 & 110 & 1,26 \\
\cellcolor{celeste}\bfseries Largueros CF + bisagras & 11,0 & 103 & 1,13 \\
\cellcolor{celeste}\bfseries Pivotes: pernos M4, resortes, bloques & 8,0 & 14 & 0,11 \\
\cellcolor{celeste}\bfseries Cápsula de nicrom (PTFE) & 5,0 & 194 & 0,97 \\
\cellcolor{celeste}\bfseries Paneles solares (2) & 9,0 & 160 & 1,44 \\
\cellcolor{celeste}\bfseries Baliza + batería propia & 15,0 & 12 & 0,18 \\
\cellcolor{celeste}\bfseries Cámara de nadir & 11,5 & 10 & 0,12 \\
\cellcolor{celeste}\bfseries Cámara de cenit & 11,5 & 203 & 2,33 \\
\cellcolor{celeste}\bfseries Cableado y conectores & 20,0 & 120 & 2,40 \\
\cellcolor{celeste}\bfseries Pack 2 $\times$ 18650 & 92,0 & 60 & 5,52 \\
\cellcolor{celeste}\bfseries Cuna en U (PETG) & 10,0 & 60 & 0,60 \\
\cellcolor{celeste}\bfseries M2 & 13,2 & 85 & 1,12 \\
\cellcolor{celeste}\bfseries Tuercas de M2 (8, Al 7075) & 1,6 & 85 & 0,14 \\
\cellcolor{celeste}\bfseries PCB principal & 28,0 & 132 & 3,70 \\
\cellcolor{celeste}\bfseries Radio 915 MHz + antena & 14,5 & 128 & 1,86 \\
\cellcolor{celeste}\bfseries GNSS + antena & 12,0 & 157 & 1,88 \\
\cellcolor{celeste}\bfseries Rieles de las placas & 8,0 & 134 & 1,07 \\*
\cellcolor{celeste}\bfseries M3 & 13,2 & 175 & 2,31 \\*
\cellcolor{celeste}\bfseries Tuercas de M3 (8, Al 7075) & 1,6 & 175 & 0,28 \\*
\cellcolor{celeste}\bfseries \bfseries Total (posición nominal) & \cellcolor{celeste}\bfseries \bfseries 531,8 & \cellcolor{celeste}\bfseries \bfseries 110,7 & \cellcolor{celeste}\bfseries  \\
\end{longtable}}
\par\medskip
\nota{Masas del documento v2, con la carena recalculada (163\,g para 240\,mm con el canal y las ranuras) y las tuercas de las estaciones. $x$ desde el fondo de nadir. Los módulos de batería y de aviónica están en su posición nominal.}
\par{\small\noindent\begin{tabularx}{\linewidth}{LCCCC}
\rowcolor{azul}\color{white}\bfseries Posición & \color{white}\bfseries $x$ de M2 & \color{white}\bfseries $x$ de M3 & \color{white}\bfseries CG [mm] & \color{white}\bfseries Margen estático \\
\cellcolor{celeste}\bfseries Nominal & 83 & 173 & 110,7 & 4--16\,\% \\
\cellcolor{celeste}\bfseries Bateria abajo & 66 & 173 & 107,0 & 6--18\,\% \\
\cellcolor{celeste}\bfseries Bateria arriba & 90 & 173 & 112,2 & 3--15\,\% \\
\cellcolor{celeste}\bfseries Todo abajo & 66 & 156 & 104,5 & 7--19\,\% \\
\cellcolor{celeste}\bfseries Todo arriba & 90 & 177 & 112,8 & 3--15\,\% \\
\end{tabularx}}
\par\medskip
\fig{0.84}{v2BR_cg.png}{CG según el corrimiento de los módulos. Correr también la aviónica casi duplica el efecto. La línea roja es la estimación baja del punto neutro de la vela: el CG tiene que quedar debajo.}
\fig{1.0}{v2BR_posiciones.png}{Las dos posiciones extremas y la nominal, con el CG marcado.}
El margen estático se calcula con el punto neutro entre el 55 y el 67\,\% de la quilla (190\,mm) desde el ápice en
$x = 14$\,mm, es decir $x_\mathrm{PN}$ = 119--141\,mm, y
$\mathrm{ME} = (x_\mathrm{PN} - x_\mathrm{CG})/c_r$.

\pend{Lo que muestra la cuenta:} v2 pide el CG entre 95 y 105\,mm, pero con las masas recalculadas el nominal queda en
110,7\,mm. Solo con los dos módulos hacia nadir se llega a 104,5\,mm. El
margen estático es positivo en todo el rango, pero es chico con la estimación baja del punto neutro. Correr los módulos
hacia nadir siempre lo aumenta. Para ir más abajo hace falta lastre en M1 (10\,g en $x = 2$ bajan el CG unos 2\,mm), pero
el CanSat ya está pasado de masa.

\section{Impulso de 30\,G: prueba de caída (FEM)}
\subsection{Ensayo y modelo}
La prueba de caída de la competencia ata el paracaídas a una cuerda no elástica de 61\,cm y suelta el CanSat desde la altura
del punto de anclaje. Al tensarse la cuerda, el CanSat recibe unos 30\,G a lo largo del eje. Se verifica el paracaídas, su
anclaje y los soportes de los equipos y de la batería.
\begin{itemize}
\item \textbf{Apoyo:} el CanSat cuelga de la brida. En la opción B los anclajes son los extremos de las cuatro varillas,
fijos en las tres direcciones; en la A, cuatro cáncamos M3 roscados en M4.
\item \textbf{Carga:} inercia de 1\,G hacia nadir, escalada al pulso. Masa de 524\,g (el
CanSat sin el paracaídas).
\item \textbf{Modelo:} tetraedros cuadráticos (CalculiX). Los equipos son masa repartida en la pieza que los sostiene; en la
opción regulable el pack de 92\,g va en la cuna de M2 y la aviónica en M3.
\item \textbf{Mamparos y tuercas:} el mamparo apretado entre dos tuercas se modela pegado a la varilla.
\end{itemize}

\fig{1.0}{v2BR_fem_modelo.png}{Modelo FEM de la opción B regulable: malla por material, anclajes de la brida y carga. Derecha: el pulso de la cuerda y el factor dinámico según su duración.}
\subsection{Pulso y factor dinámico}
\begin{ecu}
\begin{equation}
v = \sqrt{2\,g\,h} = \sqrt{2\cdot 9{,}81\cdot 0{,}61} = 3{,}46~\mathrm{m/s},\qquad
T = \frac{\pi\,v}{a_\mathrm{máx}} = \frac{\pi\cdot 3{,}46}{30\cdot 9{,}81} = 37~\mathrm{ms}
\end{equation}
\end{ecu}
Se supone que la cuerda frena y devuelve al CanSat como un resorte (semiseno con rebote). Con el primer modo axial
por encima de 492\,Hz (ninguno de los ocho primeros modos es axial), el pulso es mucho más largo que el período propio y el factor dinámico queda en
1,01: el choque es casi estático. Para un pulso corto, el peor caso de cualquier duración, el factor
sube a 1,72.

\subsection{Criterio de rotura}
Sin factor de seguridad. En las piezas impresas manda la menor entre la resistencia en el plano y la tracción entre capas.

\par{\small\noindent\begin{tabularx}{\linewidth}{P{3.4cm}P{3.6cm}L}
\rowcolor{azul}\color{white}\bfseries Material & \color{white}\bfseries Pieza & \color{white}\bfseries Resistencia usada \\
\cellcolor{celeste}\bfseries PETG impreso & Carena / monocasco & 45\,MPa (von Mises); 25\,MPa de tracción entre capas \\
\cellcolor{celeste}\bfseries PETG, 40\,\% de relleno & Mamparos y cuna & 20\,MPa; 11\,MPa entre capas \\
\cellcolor{celeste}\bfseries Al 6061-T6 & Varillas M3 & 310\,MPa \\
\cellcolor{celeste}\bfseries Acero A2-70 & Cáncamos (opción A) & 700\,MPa \\
\cellcolor{celeste}\bfseries Adhesivo sobre PETG & Pegados (opción A) & 4\,MPa de corte (rango 2--6\,MPa) \\
\end{tabularx}}
\par\medskip
\subsection{Resultados}
\par{\small\noindent\begin{tabularx}{\linewidth}{P{4.0cm}CCC}
\rowcolor{azul}\color{white}\bfseries Pieza & \color{white}\bfseries A: monocasco & \color{white}\bfseries B: esqueleto & \cellcolor{rojo}\color{white}\bfseries B regulable \\
\cellcolor{celeste}\bfseries Pared de PETG & 554 (330) & 2.573 (1.512) & \cellcolor{rosa}2.471 (1.451) \\
\cellcolor{celeste}\bfseries Mamparo M1 & 1.249 (744) & 592 (348) & \cellcolor{rosa}793 (466) \\
\cellcolor{celeste}\bfseries Mamparo M2 & 1.592 (949) & 524 (308) & \cellcolor{rosa}427 (250) \\
\cellcolor{celeste}\bfseries Mamparo M3 & 3.373 (2.010) & 606 (356) & \cellcolor{rosa}530 (311) \\
\cellcolor{celeste}\bfseries Mamparo M4 (brida) & 131 (78) & 431 (253) & \cellcolor{rosa}430 (252) \\
\cellcolor{celeste}\bfseries Pegado mamparo-pared & 92 (55) & --- & \cellcolor{rosa}--- \\
\cellcolor{celeste}\bfseries Varillas o cáncamos & 505 (301) & 842 (495) & \cellcolor{rosa}752 (441) \\
\cellcolor{celeste}\bfseries Cuna de la batería & --- & --- & \cellcolor{rosa}658 (387) \\
\end{tabularx}}
\par\medskip
\nota{Pico del pulso, en G, que lleva la pieza a la rotura (percentil 99,9 de la pieza) con el pulso de la cuerda; entre paréntesis, con el pulso más desfavorable de cualquier duración. Se excluyen los elementos a menos de 4\,mm de los anclajes (introducción de la carga).}
\fig{0.86}{v2_caida_rotura_3.png}{G de rotura de cada pieza en las tres variantes. La línea punteada es la prueba de caída.}
\fig{1.0}{v2_caida_mapas_3.png}{Fracción de la rotura con 30\,G en las tres variantes (corte por el plano medio).}
\pend{Lectura:} las tres variantes pasan los 30\,G. En la opción A todo cuelga del pegado de M4, que rompe a
92\,G y depende de cuánto agarra el adhesivo sobre PETG. En B la carga va por las varillas
y el eslabón débil pasa a ser un mamparo, a 431\,G. En la opción regulable el pack cuelga de
M2: ese mamparo trabaja más, pero el mínimo sigue en 427\,G (mamparo M2),
más de 14 veces la carga del ensayo.

\section{Pendientes}
\par{\footnotesize\noindent\begin{tabularx}{\linewidth}{P{3.0cm}LL}
\rowcolor{azul}\color{white}\bfseries Pendiente & \color{white}\bfseries Por qué & \color{white}\bfseries Cómo se cierra \\
\cellcolor{celeste}\bfseries Masa (S1) & 531,8\,g recalculados, sobre el máximo de 510\,g & Bajar el relleno de los mamparos, acortar la carena, revisar la electrónica \\
\cellcolor{celeste}\bfseries Ensayo de la lijadora & Con la carena de PETG de 2\,mm el primer modo lateral queda dentro de 0--233\,Hz (MS $-0{,}55$) & Carena laminada (memoria v3) o rigidizar la carena \\
\cellcolor{celeste}\bfseries Área de la vela & Con la quilla de 190\,mm y los largueros de 225\,mm a 52\textdegree{} el área es 0,045\,m², no 0,067 & Redimensionar la vela o verificar la definición de área de v2 \\
\cellcolor{celeste}\bfseries CG en el rango de v2 & Nominal 110,7\,mm; con los módulos hacia nadir, 104,5\,mm & Medir el CG real y ajustar; ensayos de planeo \\
\cellcolor{celeste}\bfseries Tuercas que se aflojan & Vibración de la lijadora y del vuelo & Nyloc + marca de pintura; inspección después de cada ensayo \\
\cellcolor{celeste}\bfseries Prueba térmica (60\,\textdegree C) & El PETG pierde rigidez cerca de 60\,\textdegree C & Ensayo de 2\,h con la estructura integrada \\
\cellcolor{celeste}\bfseries Choque lateral & Solo se verificó la caída axial & Caso de 30\,G lateral (aterrizaje) con el mismo modelo \\
\end{tabularx}}
\par\medskip
\end{document}
